\documentclass{article}
\usepackage[utf8]{inputenc}
\usepackage{geometry}
\usepackage{enumitem}
\usepackage{hyperref}
\usepackage{xcolor}
\usepackage{graphicx}
\usepackage{array}
\usepackage{multirow}
\usepackage{amsmath, amssymb}
\usepackage{chemfig}
\usepackage{mhchem}
\usepackage{tikz}
\usepackage{setspace}
\usepackage{soul}
\usepackage{mathtools}
\usepackage{booktabs}
\usepackage{chemformula}
\usepackage{siunitx}
\usetikzlibrary{positioning, calc}
\usepackage{longtable}
\geometry{margin=1in}
\setlength{\parskip}{0.5em}
\usepackage{cancel}

\begin{document}

\begin{center}
    {\LARGE\textbf{SI Session Planning Form}}
\end{center}

\
\noindent
\begin{flushleft}
\textbf{SI Leader:} Brandon Cobb\\
\textbf{Session Week:} 15\\
\textbf{Session\#:} 15\\
\textbf{Instructor:} Professor Tramontozzi\\
\textbf{Old Plans:}\href{https://macombedu.sharepoint.com/:b:/s/ASC-SupplementalInstructionEmbeddedTutoring/EeSYqCwT_KhMolgJvXbVEjgBOyF_DevvS8BgCm1zPkQQJg?e=QJTt4o}{SharePoint}\\
\textbf{Session Goal:} Students will review and apply concepts of concentration and solution chemistry, including calculating molarity, mass percent, volume percent, and mass/volume percent. They will apply dilution principles and use solubility rules to determine whether a solution is unsaturated, saturated, or supersaturated, and to predict precipitation. Students will analyze factors affecting solubility, including temperature and the polarity of solute and solvent. They will also explore colligative properties such as freezing point depression and boiling point elevation for electrolytes and nonelectrolytes. Additionally, students will practice stoichiometric calculations, identify limiting reactants and theoretical yields, and classify chemical reactions including synthesis, decomposition, single and double replacement, combustion, neutralization, and redox reactions. By the end of this unit, students will have mastered understanding the difference between dissociation and ionization, calculating pH, pOH, \([H_3O^+]\), and \([OH^-]\), determining percent ionization for weak acids, interpreting \(K_a\) and pKa values, identifying conjugate acid-base pairs, analyzing diprotic and triprotic acids, and predicting how salts affect the pH of a solution through hydrolysis. They will also be able to apply these concepts to assess whether solutions are acidic, basic, or neutral, and to check the consistency of their calculations with proper units and significant figures.

\end{flushleft}

\vspace{1em}
\section*{\underline{Opener}}
{\centering\large\textbf{\hl{Take Attendance}}\par}\vspace{-0.25\baselineskip}
\noindent
\begin{tabular}{|p{4.2cm}|p{9.8cm}|}
\hline
\textbf{Collaborative Learning Techniques} & Turn to a Partner\\
\hline
\textbf{Learning Strategy} & Ice Breaker (Two truths and a Lie) \\
\hline
\textbf{Time} & 11:30---11:35\\
\hline
\textbf{Materials} & Notes, homework, whiteboard, markers, ready to speak\\
\hline
\textbf{Lecture/Textbook/Old Plans/Slide References} & Class roster \\
\hline
\end{tabular}

\vspace{0.5cm}
{\large\textbf{Definitions}}
\\
\begin{tabular}{|p{0.9\linewidth}|}
\hline
\begin{minipage}[t]{0.9\linewidth}
\begin{enumerate}
    \item\textbf{Turn to a Partner:}  Group members work with a partner on an assignment or discussion topic. This technique works best with group participants who have already been provided with enough background on a subject that they can immediately move to a discussion with their partner without previewing or reviewing concepts. 
   \item\textbf{Ice Breaker:} Each student states two true facts and one false statement related to chemistry concepts (or general academic experiences), and the group guesses which statement is false. This activity encourages peer interaction, sparks discussion, and lightly reviews content in a fun, low-pressure way. 
\end{enumerate}
\end{minipage} \\
\hline
\end{tabular}
\newpage

\noindent
{\large\textbf{Checking for Understanding (questions))}}\\

\begin{tabular}{|p{0.9\linewidth}|}
\hline
\begin{itemize}
    \item What is the primary goal of this activity?

    
    \item How do students know if their solution at a station is correct?
.
    
    \item What should students do if they are unsure about a question?

    
    \item How does moving between stations benefit student learning?

    
    \item What is the importance of showing work on the team slip?

    
    \item How are bonus points awarded during the activity?
n.
    
    \item What is the role of the facilitator at each station?

    
    \item How should teams handle disagreements while solving a question?

    
    \item Why is it important that only verified answers are recorded on the team slip?

    
    \item What is the wrap-up discussion intended to accomplish?
\end{itemize} \\
\hline
\end{tabular}




\section*{\underline{Activity 1}}
\noindent
\begin{tabular}{|p{4.2cm}|p{9.8cm}|}
\hline
\textbf{Collaborative Learning Techniques} & Group Discussion\\
\hline
\textbf{Learning Strategy} & Lecture Review\\
\hline
\textbf{Time} & 11:40---12:05\\
\hline
\textbf{Materials} & Notes, homework, whiteboard, markers, ready to speak, periodic table \\
\hline
\textbf{Lecture/Textbook/Old Plans/Slide References} & Chapter 
8: Solutions, Chapter 9: Chemical Reactions\\
\hline
\end{tabular}

\vspace{0.5cm}
{\large\textbf{Definitions}}
\\
\begin{tabular}{|p{0.9\linewidth}|}
\hline
\begin{minipage}[t]{0.9\linewidth}
\begin{enumerate}
    \item\textbf{Group Discussion:} One person in the group is asked to present on a topic or review material for the group and then lead the discussion for the group. This person should not be the regular group leader.
   \item\textbf{Lecture Review:} As a group summarize the lecture from the previous class. You may have to provide prompts for the students. For example, "The first concept discussed was Civil Liberties and Public Policy, what did the profressor highlight regarding this?" You may want to ask them to try summarizing without looking at their notes; however, if they are having a difficult time remembering, tell them to refer to their notes.
\end{enumerate}
\end{minipage} \\
\hline
\end{tabular}

\newpage
\noindent
\begin{tabular}{|p{0.9\linewidth}|}
\hline
{\large\textbf{Activity Overview}} \\
\begin{enumerate}
    \item \textbf{Purpose:} Provide students with a dynamic, interactive review session that encourages collaboration, reasoning, and problem-solving across multiple concept areas. The activity simulates a “challenge pathway” in which teams work through multiple stations to reinforce mastery of previously taught material.
    
    \item \textbf{Structure:} 
    \begin{itemize}
        \item Students work in pairs or small teams.
        \item The classroom is organized into multiple stations, each focused on a specific conceptual area (e.g., calculations, prediction exercises, conceptual reasoning).
        \item Each station contains laminated question cards with problems, a space for showing work, and instructions for submitting answers.
        \item Teams choose a station to start and move freely between stations after completing questions.
    \end{itemize}
    
    \item \textbf{Materials:} 
    \begin{itemize}
        \item Laminated question cards for each station.
        \item Team slips or half-sheets for recording verified answers.
        \item Answer keys for instructors to verify correctness.
        \item Optional props or visual aids to enhance engagement (charts, diagrams, or manipulatives).
    \end{itemize}
\end{enumerate} \\
\hline
\end{tabular}
\newpage

\noindent
{\large \textbf{Instructions (detailed activity breakdown below) continued...}} \\
\vspace{0.5cm}
\noindent
\begin{tabular}{|p{0.9\linewidth}|}
\hline
{\large\textbf{Facilitation Instructions}} \\
\begin{enumerate}
    \item \textbf{Group Management:} 
    \begin{itemize}
        \item Assign students to pairs or small teams.
        \item Allow teams to start at any station to reduce bottlenecks and encourage autonomy.
    \end{itemize}

    \item \textbf{Station Interaction:} 
    \begin{itemize}
        \item Students select a question card, collaborate to solve it, and record their answer on their team slip.
        \item Encourage students to show their reasoning clearly; provide hints or clarifying questions but avoid giving answers.
        \item Verify each answer promptly before students move to a new station.
        \item Teams may revisit stations or select a new one once a problem is verified.
    \end{itemize}

    \item \textbf{Instructor Role:}
    \begin{itemize}
        \item Circulate continuously, monitoring collaboration and engagement.
        \item Ask guiding questions to promote deeper thinking or correct misconceptions.
        \item Ensure teams follow the rules and document answers properly.
        \item Keep the activity flowing by encouraging timely progress from station to station.
    \end{itemize}
\end{enumerate} \\
\hline
\end{tabular}

\vspace{0.5cm}
\noindent
\begin{tabular}{|p{0.9\linewidth}|}
\hline
{\large\textbf{Rules and Expectations}} \\
\begin{enumerate}
    \item Teams must collaborate fully; no switching partners mid-activity.
    \item All work must be shown and explained before an answer is verified.
    \item Direct copying of answers from other teams is prohibited.
    \item Teams should communicate their reasoning with peers and instructors, but hints should remain guiding rather than directive.
    \item Verified answers are the only ones that count toward scoring.
\end{enumerate} \\
\hline
\end{tabular}
\newpage

\noindent
{\large \textbf{Instructions (detailed activity breakdown below) continued...}} \\
\vspace{0.5cm}
\noindent
\begin{tabular}{|p{0.9\linewidth}|}
\hline
{\large\textbf{Scoring and Feedback}} \\
\begin{enumerate}
    \item Each correctly verified answer earns points for the team.
    \item Bonus points may be awarded for:
    \begin{itemize}
        \item Clear demonstration of reasoning.
        \item Correct use of units and notation.
        \item Effective collaboration and communication.
    \end{itemize}
    \item Maintain a visible leaderboard to promote friendly competition and motivation.
    \item Provide immediate verbal feedback at each station to reinforce correct approaches.
\end{enumerate} \\
\hline
\end{tabular}

\vspace{0.5cm}
\noindent
\begin{tabular}{|p{0.9\linewidth}|}
\hline
{\large\textbf{Wrap-Up and Reflection}} \\
\begin{enumerate}
    \item After completion, collect all team slips and review common challenges observed during the activity.
    \item Highlight effective problem-solving strategies and address common misconceptions.
    \item Facilitate a brief discussion prompting students to reflect on:
    \begin{itemize}
        \item Which types of questions were easiest or most challenging.
        \item How collaboration improved understanding.
        \item Key takeaways for upcoming assessments.
    \end{itemize}
    \item Summarize the conceptual links between all station topics and reinforce the practical application of the skills practiced.
\end{enumerate} \\
\hline
\end{tabular}
\newpage
\noindent{\large\textbf{Content (fully worked out problems \& answers)}}\\[0.2em]
\vspace{0.2cm}

\noindent\begin{tabular}{|p{0.9\linewidth}|}
\hline
\textbf{(1)} A Brönsted-Lowry acid is defined as a substance that: \\[0.2cm]
\begin{minipage}[t]{0.85\linewidth}
\begin{enumerate}[label=\alph*.]
\item Donates a proton
\item Accepts a proton
\item Releases electrons
\item Accepts electrons
\end{enumerate}
\end{minipage}\\[0.2cm]
\hline
\end{tabular}

\noindent\begin{tabular}{|p{0.9\linewidth}|}
\hline
\textbf{(2)} What is the oxidation number of chlorine in NaCl?\\[0.2cm]
\begin{minipage}[t]{0.85\linewidth}
\begin{enumerate}[label=\alph*.]
\item +1
\item -1
\item 0
\item +2
\end{enumerate}
\end{minipage}\\[0.2cm]
\hline
\end{tabular}

\noindent\begin{tabular}{|p{0.9\linewidth}|}
\hline
\textbf{(3)} Which of the following reactions is a redox reaction?\\[0.2cm]
\begin{minipage}[t]{0.85\linewidth}
\begin{enumerate}[label=\alph*.]
\item Zn + CuSO$_4$ $\rightarrow$ ZnSO$_4$ + Cu\\[0.1cm]
\item AgNO$_3$ + NaCl $\rightarrow$ AgCl + NaNO$_3$\\[0.1cm]
\item CaCO$_3$ $\rightarrow$ CaO + CO$_2$\\[0.1cm]
\item H$_2$O + SO$_3$ $\rightarrow$ H$_2$SO$_4$\\[0.1cm]
\end{enumerate}
\end{minipage}\\[0.2cm]
\hline
\end{tabular}

\noindent\begin{tabular}{|p{0.9\linewidth}|}
\hline
\textbf{(4)} In the reaction $\text{H}_2\text{PO}_4^- \rightleftharpoons \text{HPO}_4^{2-} + \text{H}^+$, the conjugate base is: \\[0.2cm]
\begin{minipage}[t]{0.85\linewidth}
\begin{enumerate}[label=\alph*.]
\item H$_2$PO$_4^-$
\item HPO$_4^{2-}$
\item H$^+$
\item None of the above
\end{enumerate}
\end{minipage}\\[0.2cm]
\hline
\end{tabular}
\newpage

\noindent\begin{tabular}{|p{0.9\linewidth}|}
\hline
\textbf{(5)} For the equilibrium $\text{C}(s) + \text{CO}_2(g) \rightleftharpoons 2\text{CO}(g)$, what happens when pressure increases?\\[0.2cm]
\begin{minipage}[t]{0.85\linewidth}
\begin{enumerate}[label=\alph*.]
\item The system shifts left because fewer gas moles are on the CO$_2$ side
\item The system shifts right because more gas moles are produced
\item No effect because solids and gases respond equally
\item Pressure changes only affect liquids, not gases
\end{enumerate}
\end{minipage}\\[0.2cm]
\hline
\end{tabular}


\noindent\begin{tabular}{|p{0.9\linewidth}|}
\hline
\textbf{(6)} What is the oxidation number of sulfur in SO$_2$?\\[0.2cm]
\begin{minipage}[t]{0.85\linewidth}
\begin{enumerate}[label=\alph*.]
\item +2
\item +4
\item +6
\item 0
\end{enumerate}
\end{minipage}\\[0.2cm]
\hline
\end{tabular}

\noindent\begin{tabular}{|p{0.9\linewidth}|}
\hline
\textbf{(7)} In the reaction $\text{HCl} + \text{H}_2\text{O} \rightarrow \text{H}_3\text{O}^+ + \text{Cl}^-$, which is the Brönsted-Lowry acid?\\[0.2cm]
\begin{minipage}[t]{0.85\linewidth}
\begin{enumerate}[label=\alph*.]
\item HCl
\item H$_2$O
\item H$_3$O$^+$
\item Cl$^-$
\end{enumerate}
\end{minipage}\\[0.2cm]
\hline
\end{tabular}

\noindent\begin{tabular}{|p{0.9\linewidth}|}
\hline
\textbf{(8)} For the exothermic reaction $2\text{SO}_2 + \text{O}_2 \rightarrow 2\text{SO}_3$, increasing temperature will: \\[0.2cm]
\begin{minipage}[t]{0.85\linewidth}
\begin{enumerate}[label=\alph*.]
\item Increase SO$_3$ production
\item Shift equilibrium to left
\item Not affect equilibrium
\item Increase both forward and reverse rate equally
\end{enumerate}
\end{minipage}\\[0.2cm]
\hline
\end{tabular}
\newpage

\noindent\begin{tabular}{|p{0.9\linewidth}|}
\hline
\textbf{(9)} Calcium hydroxide reacts with hydrochloric acid to form calcium chloride and water. Which statement best describes the acid-base behavior occurring in this reaction?\\[0.2cm]
Ca(OH)$_2$ + 2HCl $\rightarrow$ CaCl$_2$ + 2H$_2$O\\[0.2cm]
\begin{minipage}[t]{0.85\linewidth}
\begin{enumerate}[label=\alph*.]
\item HCl donates protons to OH$^-$ ions, forming water as Ca$^{2+}$ pairs with Cl$^-$
\item Ca(OH)$_2$ donates protons to HCl, forming hydrogen gas
\item Both reactants accept protons, preventing water formation
\item Ca$^{2+}$ acts as the proton donor while Cl$^-$ acts as the proton acceptor
\end{enumerate}
\end{minipage}\\[0.2cm]
\hline
\end{tabular}

\noindent\begin{tabular}{|p{0.9\linewidth}|}
\hline
\textbf{(10)} For the reaction $\text{R} + \text{S} \rightleftharpoons \text{RS}$, which change could temporarily increase Q?\\[0.2cm]
\begin{minipage}[t]{0.85\linewidth}
\begin{enumerate}[label=\alph*.]
\item Adding RS
\item Adding R
\item Removing S
\item Diluting all species equally
\end{enumerate}
\end{minipage}\\[0.2cm]
\hline
\end{tabular}

\noindent\begin{tabular}{|p{0.9\linewidth}|}
\hline
\textbf{(11)} For the reaction $2\text{A} \rightleftharpoons \text{B}$, which expression is correct?\\[0.2cm]
\begin{minipage}[t]{0.85\linewidth}
\begin{enumerate}[label=\alph*.]
\item $K_{eq}=\frac{[\text{B}]}{[\text{A}]^2}$
\item $K_{eq}=\frac{[\text{A}]^2}{[\text{B}]}$
\item $K_{eq}=\frac{1}{[\text{A}][\text{B}]}$
\item $K_{eq}=[\text{A}][\text{B}]$
\end{enumerate}
\end{minipage}\\[0.2cm]
\hline
\end{tabular}

\noindent\begin{tabular}{|p{0.9\linewidth}|}
\hline
\textbf{(12)} Identify the reaction type:\\[0.2cm]
BaCl$_2$ + Na$_2$SO$_4$ $\rightarrow$ BaSO$_4$ + 2NaCl\\[0.2cm]
\begin{minipage}[t]{0.85\linewidth}
\begin{enumerate}[label=\alph*.]
\item Combination
\item Decomposition
\item Exchange
\item Single replacement
\end{enumerate}
\end{minipage}\\[0.2cm]
\hline
\end{tabular}
\newpage

\noindent\begin{tabular}{|p{0.9\linewidth}|}
\hline
\textbf{(13)} In a redox reaction, the oxidizing agent is the species that:\\[0.2cm]
\begin{minipage}[t]{0.85\linewidth}
\begin{enumerate}[label=\alph*.]
\item Loses electrons
\item Gains electrons
\item Is always neutral
\item Has the highest atomic mass
\end{enumerate}
\end{minipage}\\[0.2cm]
\hline
\end{tabular}

\noindent\begin{tabular}{|p{0.9\linewidth}|}
\hline
\textbf{(14)} For the endothermic reaction $\text{N}_2\text{O}_4 \rightarrow 2\text{NO}_2$, decreasing temperature favors: \\[0.2cm]
\begin{minipage}[t]{0.85\linewidth}
\begin{enumerate}[label=\alph*.]
\item Formation of NO$_2$
\item Formation of N$_2$O$_4$
\item Equal amounts of each
\item Faster formation of NO$_2$ only
\end{enumerate}
\end{minipage}\\[0.2cm]
\hline
\end{tabular}

\noindent\begin{tabular}{|p{0.9\linewidth}|}
\hline
\textbf{(15)} For the equilibrium $ \text{PCl}_5(g) \rightleftharpoons \text{PCl}_3(g) + \text{Cl}_2(g)$, increasing pressure causes the reaction to shift: \\[0.2cm]
\begin{minipage}[t]{0.85\linewidth}
\begin{enumerate}[label=\alph*.]
\item Left, toward PCl$_5$
\item Right, toward PCl$_3$ and Cl$_2$
\item No shift because gases are unaffected by pressure
\item Toward whichever side has the higher total mass
\end{enumerate}
\end{minipage}\\[0.2cm]
\hline
\end{tabular}


\noindent\begin{tabular}{|p{0.9\linewidth}|}
\hline
\textbf{(16)} In the reaction 2Na + Cl$_2$ $\rightarrow$ 2NaCl, which is the oxidizing agent?\\[0.2cm]
\begin{minipage}[t]{0.85\linewidth}
\begin{enumerate}[label=\alph*.]
\item Na
\item Cl$_2$
\item NaCl
\item Both Na and Cl$_2$
\end{enumerate}
\end{minipage}\\[0.2cm]
\hline
\end{tabular}
\newpage

\noindent\begin{tabular}{|p{0.9\linewidth}|}
\hline
\textbf{(17)} For the reaction $\text{A} + \text{C} \rightleftharpoons \text{AC}$ with $K_{eq} = 0.02$, which is true?\\[0.2cm]
\begin{minipage}[t]{0.85\linewidth}
\begin{enumerate}[label=\alph*.]
\item Reactants are strongly favored
\item Products are strongly favored
\item Equal amounts of A and AC form
\item Keq must be recalculated using solids only
\end{enumerate}
\end{minipage}\\[0.2cm]
\hline
\end{tabular}

\noindent\begin{tabular}{|p{0.9\linewidth}|}
\hline
\textbf{(18)} According to the Arrhenius definition, which statement correctly describes an acid and a base?\\[0.2cm]
\begin{minipage}[t]{0.85\linewidth}
\begin{enumerate}[label=\alph*.]
\item An acid produces H$^+$ in water, and a base produces OH$^-$ in water
\item An acid donates a proton, and a base accepts a proton
\item An acid accepts an electron pair, and a base donates an electron pair
\item An acid and base both increase the concentration of ions equally
\end{enumerate}
\end{minipage}\\[0.2cm]
\hline
\end{tabular}


\noindent\begin{tabular}{|p{0.9\linewidth}|}
\hline
\textbf{(19)} What type of reaction occurs when methane burns in oxygen to produce carbon dioxide and water?\\[0.2cm]
CH$_4$ + 2O$_2$ $\rightarrow$ CO$_2$ + 2H$_2$O\\[0.2cm]
\begin{minipage}[t]{0.85\linewidth}
\begin{enumerate}[label=\alph*.]
\item Combination
\item Decomposition
\item Combustion
\item Single displacement
\end{enumerate}
\end{minipage}\\[0.2cm]
\hline
\end{tabular}

\noindent\begin{tabular}{|p{0.9\linewidth}|}
\hline
\textbf{(20)} What type of chemical reaction occurs between sodium phosphate and calcium nitrate to form calcium phosphate and sodium nitrate?\\[0.2cm]
2Na$_3$PO$_4$ + 3Ca(NO$_3$)$_2$ $\rightarrow$ Ca$_3$(PO$_4$)$_2$ + 6NaNO$_3$\\[0.2cm]
\begin{minipage}[t]{0.85\linewidth}
\begin{enumerate}[label=\alph*.]
\item Decomposition
\item Exchange
\item Combination
\item Combustion
\end{enumerate}
\end{minipage}\\[0.2cm]
\hline
\end{tabular}
\newpage

\noindent\begin{tabular}{|p{0.9\linewidth}|}
\hline
\textbf{(21)} In the reaction $\text{CO}(g) + 2\text{H}_2(g) \rightleftharpoons \text{CH}_3\text{OH}(g)$, which change increases methanol production? \\[0.2cm]
\begin{minipage}[t]{0.85\linewidth}
\begin{enumerate}[label=\alph*.]
\item Increasing pressure because the product side has fewer gas moles
\item Decreasing pressure because the reactant side has more gas moles
\item Increasing volume of the container
\item Adding more inert gas at constant volume
\end{enumerate}
\end{minipage}\\[0.2cm]
\hline
\end{tabular}


\noindent\begin{tabular}{|p{0.9\linewidth}|}
\hline
\textbf{(22)} If $Q > K_{eq}$ for the reaction $\text{F} \rightleftharpoons \text{G}$, which direction will the system shift?\\[0.2cm]
\begin{minipage}[t]{0.85\linewidth}
\begin{enumerate}[label=\alph*.]
\item Toward reactants
\item Toward products
\item No shift
\item The reaction stops
\end{enumerate}
\end{minipage}\\[0.2cm]
\hline
\end{tabular}

\noindent\begin{tabular}{|p{0.9\linewidth}|}
\hline
\textbf{(23)} When sodium reacts with chlorine to produce sodium chloride, what type of reaction is this?\\[0.2cm]
2Na + Cl$_2$ $\rightarrow$ 2NaCl\\[0.2cm]
\begin{minipage}[t]{0.85\linewidth}
\begin{enumerate}[label=\alph*.]
\item Decomposition
\item Combustion
\item Exchange
\item Combination
\end{enumerate}
\end{minipage}\\[0.2cm]
\hline
\end{tabular}

\noindent\begin{tabular}{|p{0.9\linewidth}|}
\hline
\textbf{(24)} For the reaction $\text{P} \rightleftharpoons \text{Q}$, which condition must be true at equilibrium?\\[0.2cm]
\begin{minipage}[t]{0.85\linewidth}
\begin{enumerate}[label=\alph*.]
\item Rate forward = rate reverse
\item $[P]$ = [Q]
\item $K_{eq}$ = 0
\item All reactant is converted to product
\end{enumerate}
\end{minipage}\\[0.2cm]
\hline
\end{tabular}
\newpage

\noindent\begin{tabular}{|p{0.9\linewidth}|}
\hline
\textbf{(25)} For the reaction $\text{A} + \text{heat} \rightarrow \text{B}$, which statement is correct? \\[0.2cm]
\begin{minipage}[t]{0.85\linewidth}
\begin{enumerate}[label=\alph*.]
\item Adding heat increases [B]
\item Adding heat decreases [B]
\item Removing heat increases [B]
\item Heat has no effect on equilibrium
\end{enumerate}
\end{minipage}\\[0.2cm]
\hline
\end{tabular}

\noindent\begin{tabular}{|p{0.9\linewidth}|}
\hline
\textbf{(26)} The burning of benzene in air to form carbon dioxide and water is an example of what type of reaction?\\[0.2cm]
2C$_6$H$_6$ + 15O$_2$ $\rightarrow$ 12CO$_2$ + 6H$_2$O\\[0.2cm]
\begin{minipage}[t]{0.85\linewidth}
\begin{enumerate}[label=\alph*.]
\item Oxidation-reduction
\item Combustion
\item Decomposition
\item Neutralization
\end{enumerate}
\end{minipage}\\[0.2cm]
\hline
\end{tabular}

\noindent\begin{tabular}{|p{0.9\linewidth}|}
\hline
\textbf{(27)} Which element is being reduced in the reaction Fe$_2$O$_3$ + 3CO $\rightarrow$ 2Fe + 3CO$_2$?\\[0.2cm]
\begin{minipage}[t]{0.85\linewidth}
\begin{enumerate}[label=\alph*.]
\item Fe
\item O
\item C
\item CO$_2$
\end{enumerate}
\end{minipage}\\[0.2cm]
\hline
\end{tabular}

\noindent\begin{tabular}{|p{0.9\linewidth}|}
\hline
\textbf{(28)} What happens to an atom's oxidation number when it is oxidized?\\[0.2cm]
\begin{minipage}[t]{0.85\linewidth}
\begin{enumerate}[label=\alph*.]
\item It increases
\item It decreases
\item It stays the same
\item It becomes zero
\end{enumerate}
\end{minipage}\\[0.2cm]
\hline
\end{tabular}
\newpage

\noindent\begin{tabular}{|p{0.9\linewidth}|}
\hline
\textbf{(29)} When ammonium chloride reacts with silver nitrate to form silver chloride and ammonium nitrate, the reaction is classified as:\\[0.2cm]
NH$_4$Cl + AgNO$_3$ $\rightarrow$ AgCl + NH$_4$NO$_3$\\[0.2cm]
\begin{minipage}[t]{0.85\linewidth}
\begin{enumerate}[label=\alph*.]
\item Combination
\item Exchange
\item Decomposition
\item Combustion
\end{enumerate}
\end{minipage}\\[0.2cm]
\hline
\end{tabular}

\noindent\begin{tabular}{|p{0.9\linewidth}|}
\hline
\textbf{(30)} What type of reaction occurs when potassium iodide reacts with lead(II) nitrate to form lead(II) iodide and potassium nitrate?\\[0.2cm]
2KI + Pb(NO$_3$)$_2$ $\rightarrow$ 2KNO$_3$ + PbI$_2$\\[0.2cm]
\begin{minipage}[t]{0.85\linewidth}
\begin{enumerate}[label=\alph*.]
\item Exchange
\item Combination
\item Decomposition
\item Combustion
\end{enumerate}
\end{minipage}\\[0.2cm]
\hline
\end{tabular}

\noindent\begin{tabular}{|p{0.9\linewidth}|}
\hline
\textbf{(31)} When hydrochloric acid reacts with sodium hydroxide to form sodium chloride and water, which statement best describes the acid-base process taking place?\\[0.2cm]
HCl + NaOH $\rightarrow$ NaCl + H$_2$O\\[0.2cm]
\begin{minipage}[t]{0.85\linewidth}
\begin{enumerate}[label=\alph*.]
\item HCl donates a proton to OH$^-$, forming water as Na$^+$ pairs with Cl$^-$
\item NaOH donates protons to HCl, forming hydrogen gas
\item Both reactants act as proton donors, preventing neutralization
\item Cl$^-$ acts as the proton donor while Na$^+$ accepts the proton
\end{enumerate}
\end{minipage}\\[0.2cm]
\hline
\end{tabular}

\noindent\begin{tabular}{|p{0.9\linewidth}|}
\hline
\textbf{(32)} In the reaction $\text{NH}_3 + \text{H}_2\text{O} \rightarrow \text{NH}_4^+ + \text{OH}^-$, which species acts as the Brönsted-Lowry base?\\[0.2cm]
\begin{minipage}[t]{0.85\linewidth}
\begin{enumerate}[label=\alph*.]
\item NH$_3$
\item H$_2$O
\item NH$_4^+$
\item OH$^-$
\end{enumerate}
\end{minipage}\\[0.2cm]
\hline
\end{tabular}
\newpage

\noindent\begin{tabular}{|p{0.9\linewidth}|}
\hline
\textbf{(33)} Identify the reaction type:\\[0.2cm]
2H$_2$O$_2$ $\rightarrow$ 2H$_2$O + O$_2$\\[0.2cm]
\begin{minipage}[t]{0.85\linewidth}
\begin{enumerate}[label=\alph*.]
\item Combination
\item Combustion
\item Decomposition
\item Exchange
\end{enumerate}
\end{minipage}\\[0.2cm]
\hline
\end{tabular}

\noindent\begin{tabular}{|p{0.9\linewidth}|}
\hline
\textbf{(34)} Which pair represents a conjugate acid-base pair?\\[0.2cm]
\begin{minipage}[t]{0.85\linewidth}
\begin{enumerate}[label=\alph*.]
\item NH$_3$/NH$_4^+$
\item OH$^-$/O$^{2-}$
\item Na$^+$/Na
\item Cl$^-$/Cl$_2$
\end{enumerate}
\end{minipage}\\[0.2cm]
\hline
\end{tabular}

\noindent\begin{tabular}{|p{0.9\linewidth}|}
\hline
\textbf{(35)} Which of the following elements always has an oxidation number of +1 in compounds?\\[0.2cm]
\begin{minipage}[t]{0.85\linewidth}
\begin{enumerate}[label=\alph*.]
\item Sodium
\item Oxygen
\item Fluorine
\item Sulfur
\end{enumerate}
\end{minipage}\\[0.2cm]
\hline
\end{tabular}

\noindent\begin{tabular}{|p{0.9\linewidth}|}
\hline
\textbf{(36)} What is the oxidation number of manganese in KMnO$_4$?\\[0.2cm]
\begin{minipage}[t]{0.85\linewidth}
\begin{enumerate}[label=\alph*.]
\item +2
\item +4
\item +6
\item +7
\end{enumerate}
\end{minipage}\\[0.2cm]
\hline
\end{tabular}
\newpage

\noindent\begin{tabular}{|p{0.9\linewidth}|}
\hline
\textbf{(37)} In a redox reaction, the reducing agent is the substance that:\\[0.2cm]
\begin{minipage}[t]{0.85\linewidth}
\begin{enumerate}[label=\alph*.]
\item Gains electrons
\item Loses electrons
\item Is neither oxidized nor reduced
\item Lowers its oxidation number
\end{enumerate}
\end{minipage}\\[0.2cm]
\hline
\end{tabular}

\noindent\begin{tabular}{|p{0.9\linewidth}|}
\hline
\textbf{(38)} For the reaction $2\text{SO}_2(g) + \text{O}_2(g) \rightleftharpoons 2\text{SO}_3(g)$, decreasing pressure favors: \\[0.2cm]
\begin{minipage}[t]{0.85\linewidth}
\begin{enumerate}[label=\alph*.]
\item Formation of SO$_3$
\item Formation of SO$_2$ and O$_2$
\item No change in equilibrium
\item Formation of whichever side absorbs heat
\end{enumerate}
\end{minipage}\\[0.2cm]
\hline
\end{tabular}


\noindent\begin{tabular}{|p{0.9\linewidth}|}
\hline
\textbf{(39)} For the endothermic reaction $\text{heat} + \text{H}_2 + \text{I}_2 \rightarrow 2\text{HI}$, decreasing temperature shifts equilibrium: \\[0.2cm]
\begin{minipage}[t]{0.85\linewidth}
\begin{enumerate}[label=\alph*.]
\item Toward products
\item Toward reactants
\item No shift
\item Depends only on pressure
\end{enumerate}
\end{minipage}\\[0.2cm]
\hline
\end{tabular}

\noindent\begin{tabular}{|p{0.9\linewidth}|}
\hline
\textbf{(40)} For the reaction $\text{A} + \text{B} \rightleftharpoons \text{C}$, if $K_{eq} = 12$, which statement is true?\\[0.2cm]
\begin{minipage}[t]{0.85\linewidth}
\begin{enumerate}[label=\alph*.]
\item Products are favored
\item Reactants are favored
\item Neither side is favored
\item The reaction goes to completion
\end{enumerate}
\end{minipage}\\[0.2cm]
\hline
\end{tabular}
\newpage

\noindent\begin{tabular}{|p{0.9\linewidth}|}
\hline
\textbf{(41)} What type of reaction occurs when lithium sulfate reacts with barium nitrate to form barium sulfate and lithium nitrate?\\[0.2cm]
Li$_2$SO$_4$ + Ba(NO$_3$)$_2$ $\rightarrow$ BaSO$_4$ + 2LiNO$_3$\\[0.2cm]
\begin{minipage}[t]{0.85\linewidth}
\begin{enumerate}[label=\alph*.]
\item Combination
\item Decomposition
\item Exchange
\item Combustion
\end{enumerate}
\end{minipage}\\[0.2cm]
\hline
\end{tabular}

\noindent\begin{tabular}{|p{0.9\linewidth}|}
\hline
\textbf{(42)} For the reaction $2\text{X} \rightleftharpoons \text{Y}$, $K_{eq} = 1.0$. Which is true at equilibrium?\\[0.2cm]
\begin{minipage}[t]{0.85\linewidth}
\begin{enumerate}[label=\alph*.]
\item $[Y]$ = [X]
\item $[Y]$ = $[X]^2$
\item $[X]$ = 2[Y]
\item $K_{eq}$ cannot equal 1 for reversible reactions
\end{enumerate}
\end{minipage}\\[0.2cm]
\hline
\end{tabular}

\noindent\begin{tabular}{|p{0.9\linewidth}|}
\hline
\textbf{(43)} What kind of reaction is the burning of glucose in oxygen to form carbon dioxide and water?\\[0.2cm]
C$_6$H$_{12}$O$_6$ + 6O$_2$ $\rightarrow$  6CO$_2$ + 6H$_2$O\\[0.2cm]
\begin{minipage}[t]{0.85\linewidth}
\begin{enumerate}[label=\alph*.]
\item Double displacement
\item Combustion
\item Decomposition
\item Combination
\end{enumerate}
\end{minipage}\\[0.2cm]
\hline
\end{tabular}

\noindent\begin{tabular}{|p{0.9\linewidth}|}
\hline
\textbf{(44)} A substance that gains electrons during a reaction is:\\[0.2cm]
\begin{minipage}[t]{0.85\linewidth}
\begin{enumerate}[label=\alph*.]
\item Oxidized
\item Reduced
\item Neutralized
\item Dissociated
\end{enumerate}
\end{minipage}\\[0.2cm]
\hline
\end{tabular}
\newpage

\noindent\begin{tabular}{|p{0.9\linewidth}|}
\hline
\textbf{(45)} Which of the following represents a redox reaction?\\[0.2cm]
\begin{minipage}[t]{0.85\linewidth}
\begin{enumerate}[label=\alph*.]
\item HCl + NaOH $\rightarrow$ NaCl + H$_2$O\\[0.1cm]
\item Zn + 2HCl $\rightarrow$ ZnCl$_2$ + H$_2$\\[0.1cm]
\item CaCO$_3$ $\rightarrow$ CaO + CO$_2$\\[0.1cm]
\item NaOH + HNO$_3$ $\rightarrow$ NaNO$_3$ + H$_2$O\\[0.1cm]
\end{enumerate}

\end{minipage}\\[0.2cm]
\hline
\end{tabular}

\noindent\begin{tabular}{|p{0.9\linewidth}|}
\hline
\textbf{(46)} For the exothermic decomposition $2\text{NO}_{(g)} \rightarrow {\text{N}_{2}}_{(g)} + {\text{O}_{2}}_{(g)} + \text{heat}$, adding heat will favor: \\[0.2cm]
\begin{minipage}[t]{0.85\linewidth}
\begin{enumerate}[label=\alph*.]
\item Formation of N$_2$ and O$_2$
\item Formation of NO
\item No change
\item Whichever side has fewer gas particles
\end{enumerate}
\end{minipage}\\[0.2cm]
\hline
\end{tabular}

\noindent\begin{tabular}{|p{0.9\linewidth}|}
\hline
\textbf{(47)} What type of reaction occurs when silver nitrate reacts with sodium chloride to produce silver chloride and sodium nitrate?\\[0.2cm]
AgNO$_3$ + NaCl $\rightarrow$ AgCl + NaNO$_3$\\[0.2cm]
\begin{minipage}[t]{0.85\linewidth}
\begin{enumerate}[label=\alph*.]
\item Combination
\item Decomposition
\item Exchange
\item Combustion
\end{enumerate}
\end{minipage}\\[0.2cm]
\hline
\end{tabular}

\noindent\begin{tabular}{|p{0.9\linewidth}|}
\hline
\textbf{(48)} For the reaction $\text{H}_2 + \text{I}_2 \rightleftharpoons 2\text{HI}$, which change increases $Q$?\\[0.2cm]
\begin{minipage}[t]{0.85\linewidth}
\begin{enumerate}[label=\alph*.]
\item Adding HI
\item Removing HI
\item Adding H$_2$
\item Removing heat
\end{enumerate}
\end{minipage}\\[0.2cm]
\hline
\end{tabular}
\newpage

\noindent\begin{tabular}{|p{0.9\linewidth}|}
\hline
\textbf{(49)} What type of chemical reaction does ethanol undergo in air?\\[0.2cm]
C$_2$H$_5$OH + 3O$_2$ $\rightarrow$ 2CO$_2$ + 3H$_2$O\\[0.2cm]
\begin{minipage}[t]{0.85\linewidth}
\begin{enumerate}[label=\alph*.]
\item Neutralization
\item Single replacement
\item Combustion
\item Precipitation
\end{enumerate}
\end{minipage}\\[0.2cm]
\hline
\end{tabular}

\noindent\begin{tabular}{|p{0.9\linewidth}|}
\hline
\textbf{(50)} For the reaction $\text{A} + \text{B} + \text{heat} \rightarrow \text{C}$, which way does equilibrium shift when heat is added?\\[0.2cm]
\begin{minipage}[t]{0.85\linewidth}
\begin{enumerate}[label=\alph*.]
\item Toward products
\item Toward reactants
\item No change
\item It depends on pressure only
\end{enumerate}
\end{minipage}\\[0.2cm]
\hline
\end{tabular}

\noindent\begin{tabular}{|p{0.9\linewidth}|}
\hline
\textbf{(51)} For the exothermic reaction $\text{X} + \text{Y} \rightarrow \text{Z} + \text{heat}$, removing heat shifts equilibrium: \\[0.2cm]
\begin{minipage}[t]{0.85\linewidth}
\begin{enumerate}[label=\alph*.]
\item Toward reactants
\item Toward products
\item Causes no shift
\item Only affects rate, not equilibrium
\end{enumerate}
\end{minipage}\\[0.2cm]
\hline
\end{tabular}

\noindent\begin{tabular}{|p{0.9\linewidth}|}
\hline
\textbf{(52)} In the equilibrium $\text{HF} + \text{H}_2\text{O} \rightleftharpoons \text{H}_3\text{O}^+ + \text{F}^-$, which is the conjugate acid of F$^-$? \\[0.2cm]
\begin{minipage}[t]{0.85\linewidth}
\begin{enumerate}[label=\alph*.]
\item HF
\item H$_3$O$^+$
\item H$_2$O
\item F$^-$ has no conjugate acid
\end{enumerate}
\end{minipage}\\[0.2cm]
\hline
\end{tabular}
\newpage

\noindent\begin{tabular}{|p{0.9\linewidth}|}
\hline
\textbf{(53)} What type of reaction occurs when acetylene reacts with oxygen to produce CO$_2$ and H$_2$O?\\[0.2cm]
2C$_2$H$_2$ + 5O$_2$ $\rightarrow$ 4CO$_2$ + 2H$_2$O\\[0.2cm]
\begin{minipage}[t]{0.85\linewidth}
\begin{enumerate}[label=\alph*.]
\item Neutralization
\item Combustion
\item Combination
\item Decomposition
\end{enumerate}
\end{minipage}\\[0.2cm]
\hline
\end{tabular}

\noindent\begin{tabular}{|p{0.9\linewidth}|}
\hline
\textbf{(54)} For the reaction $\text{N}_2(g) + 3\text{H}_2(g) \rightleftharpoons 2\text{NH}_3(g)$, increasing pressure shifts the equilibrium: \\[0.2cm]
\begin{minipage}[t]{0.85\linewidth}
\begin{enumerate}[label=\alph*.]
\item Toward products because fewer gas moles are present
\item Toward reactants because more gas moles are present
\item No shift because pressure does not affect gases
\item Toward whichever side has more volume
\end{enumerate}
\end{minipage}\\[0.2cm]
\hline
\end{tabular}

\noindent\begin{tabular}{|p{0.9\linewidth}|}
\hline
\textbf{(55)} In the reaction Zn + Cu$^{2+}$ $\rightarrow$ Zn$^{2+}$ + Cu, which species is oxidized?\\[0.2cm]
\begin{minipage}[t]{0.85\linewidth}
\begin{enumerate}[label=\alph*.]
\item Cu$^{2+}$
\item Cu
\item Zn
\item H$_2$O
\end{enumerate}
\end{minipage}\\[0.2cm]
\hline
\end{tabular}

\noindent\begin{tabular}{|p{0.9\linewidth}|}
\hline
\textbf{(56)} Identify the type of reaction below:\\[0.2cm]
K$_2$SO$_4$ + Ba(NO$_3$)$_2$ $\rightarrow$ 2KNO$_3$ + BaSO$_4$\\[0.2cm]
\begin{minipage}[t]{0.85\linewidth}
\begin{enumerate}[label=\alph*.]
\item Exchange
\item Combination
\item Combustion
\item Decomposition
\end{enumerate}
\end{minipage}\\[0.2cm]
\hline
\end{tabular}
\newpage

\noindent\begin{tabular}{|p{0.9\linewidth}|}
\hline
\textbf{(57)} For the reaction $\text{A} + \text{C} \rightleftharpoons \text{AC}$ with $K_{eq} = 0.02$, which is true?\\[0.2cm]
\begin{minipage}[t]{0.85\linewidth}
\begin{enumerate}[label=\alph*.]
\item Reactants are strongly favored
\item Products are strongly favored
\item Equal amounts of A and AC form
\item $K_{eq}$ must be recalculated using solids only
\end{enumerate}
\end{minipage}\\[0.2cm]
\hline
\end{tabular}

\noindent\begin{tabular}{|p{0.9\linewidth}|}
\hline
\textbf{(58)} Which of the following compounds acts as an Arrhenius base when dissolved in water?\\[0.2cm]
\begin{minipage}[t]{0.85\linewidth}
\begin{enumerate}[label=\alph*.]
\item NaOH, because it increases the concentration of OH$^-$ in solution
\item HCl, because it donates protons to other species
\item NH$_3$, because it donates an electron pair to H$^+$
\item CO$_2$, because it produces H$^+$ in water
\end{enumerate}
\end{minipage}\\[0.2cm]
\hline
\end{tabular}


\noindent\begin{tabular}{|p{0.9\linewidth}|}
\hline
\textbf{(59)} For $\text{D} \rightleftharpoons \text{E}$, which scenario decreases the value of the reaction quotient $Q$?\\[0.2cm]
\begin{minipage}[t]{0.85\linewidth}
\begin{enumerate}[label=\alph*.]
\item Adding D
\item Removing D
\item Adding E
\item Both b. and c.
\end{enumerate}
\end{minipage}\\[0.2cm]
\hline
\end{tabular}

\noindent\begin{tabular}{|p{0.9\linewidth}|}
\hline
\textbf{(60)} For the reaction $\text{M} \rightleftharpoons \text{N}$, if $K_{eq}$ is extremely large, which is true?\\[0.2cm]
\begin{minipage}[t]{0.85\linewidth}
\begin{enumerate}[label=\alph*.]
\item Products dominate at equilibrium
\item Reactants dominate at equilibrium
\item $[M]$ must be greater than [N]
\item Q must equal zero
\end{enumerate}
\end{minipage}\\[0.2cm]
\hline
\end{tabular}
\newpage
\noindent
{\large\textbf{Checking for Understanding (questions \& answers)}}\\

\begin{tabular}{|p{0.9\linewidth}|}
\hline
\begin{itemize}
\item How does adding a nonvolatile solute affect the freezing point of a solvent?
\item Why does dissolving an ionic compound usually lower the freezing point more than a covalent one?
\item What factors determine the magnitude of freezing point depression?
\item Why does the identity (not just the amount) of the solute matter in colligative properties?
\item How does the boiling point of a solution compare to that of the pure solvent?
\item Why does adding solute increase the boiling point of a liquid?
\item What physical property remains unchanged when a solution’s freezing point is lowered?
\item Why does temperature itself not change the amount of freezing point depression?
\item Why do electrolytes and nonelectrolytes differ in their effects on colligative properties?
\item How does a neutralization reaction differ from a precipitation or redox reaction?
\item Why is acid–base neutralization classified as a double replacement reaction?
\item What distinguishes a synthesis reaction from a decomposition reaction?
\item How is a combustion reaction identified based on its reactants and products?
\item Why does every combustion reaction of a hydrocarbon produce CO$_2$ and H$_2$O?
\item How does a single replacement reaction differ from a double replacement reaction?
\item Why must the activity of elements be considered when predicting single replacement outcomes?
\item How can you identify the limiting reactant in a reaction from given amounts of reactants?
\item Why does conservation of mass require balancing chemical equations?
\item How can stoichiometric ratios be used to determine theoretical yield?
\item Why do ionic compounds dissociate in water while molecular ones may not?
\item How does solubility influence whether a precipitate forms in a double replacement reaction?
\item Why is water often called the “universal solvent” in ionic and polar systems?
\item How does polarity explain why oil and water do not mix?
\item Why do stronger intermolecular forces lead to higher boiling points?
\item How can temperature changes indicate whether a process is endothermic or exothermic?
\item Why is combustion considered an exothermic oxidation process?
\item What evidence suggests that a chemical reaction, not a physical change, has occurred?
\end{itemize}\\
\hline
\end{tabular}


\newpage

\section*{\underline{Activity 2}}
\vspace{0.2cm}
\begin{tabular}{|p{4.2cm}|p{9.8cm}|}
\hline
\textbf{Collaborative Learning Techniques} & Jigsaw\\
\hline
\textbf{Learning Strategy} & Predict Test Questions \\
\hline
\textbf{Time} & 12:05---12:30\\
\hline
\textbf{Materials} & Pencil, eraser, calculator, periodic table \\
\hline
\textbf{Lecture/Textbook/Old Plans/Slide References} & Chapter 10: Acid and Bases \\
\hline
\end{tabular}
\vspace{0.5cm}

{\large\textbf{Definition}}\\
\noindent
\begin{tabular}{|p{0.9\linewidth}|}
\hline
\begin{enumerate}
    \item\textbf{Jigsaw:} Jigsaws, when used properly, make the group as a whole dependent upon all of them in subgroups. Each group provides a peice of the puzzle. Group members are broken into smaller groups. Each small group works on some aspect of the same problem, question, or issue. They then share their part of the puzzle with the large group.
    \item\textbf{Predict Test Questions:}  Put students in groups of 2 or 3 and assign them to write a test question for a specific topic, ensuring that all topics have been convered. Ask students to write their questions on the board or on an overhead for discussion (would the professor ask the question? What is the answer? etc.). Students will have the benefit of learning to think like the teacher and they'll be able to see additional questions that other students have writen.
 
\end{enumerate} \\
\hline
\end{tabular}
\newpage

\vspace{0.5cm}
\noindent
{\large \textbf{Instructions (detailed activity breakdown below)}} \\

\begin{tabular}{|p{0.9\linewidth}|}
\hline
\begin{enumerate}
    \item At each station, place a set of laminated question cards (multiple choice and true/false). Each card includes:
        \begin{itemize}
            \item A unique question number and topic label.
            \item The possible answer choices (A–E or T/F).
            \item A mapping line showing which station each answer leads to (to simulate “pathways” through the escape room).
        \end{itemize}
        \item Provide each pair of students with a \textbf{team slip} (half-sheet) labeled with their team name and space to record up to 3 verified answers.
    \item \textbf{Facilitation:}
    \begin{itemize}
        \item Pairs may begin at \textbf{any station of their choice}.
        \item Each pair selects one question card, discusses the problem, and writes their chosen answer \textbf{on the back of the card}.
        \item Students then show their work and reasoning to the instructor or SI leader for quick verification.
        \item If the answer is correct, they record that problem on their team slip (problem number and answer) and may proceed to any other station of their choice.
        \item Pairs are encouraged to move freely between stations, aiming to correctly solve at least \textbf{three problems} during the activity period.
        \item After completing at least three problems, teams submit their slips to the instructor for scoring.
    \end{itemize}

    \item \textbf{Scoring:}
    \begin{itemize}
        \item Each correctly solved problem = \textbf{1 point}.
        \item Teams with the highest total score win or “escape” first.
        \item Bonus points may be awarded for clear reasoning or teamwork.
    \end{itemize}

    \item \textbf{Rules:}
    \begin{itemize}
        \item Teams must work collaboratively—no switching partners mid-activity.
        \item All work must be shown before a response is accepted.
        \item Only verified answers written on the team slip count toward the final score.
        \item Communication between pairs is encouraged, but direct sharing of answers is not permitted.
    \end{itemize}

    \item \textbf{Wrap-Up:}
    \begin{itemize}
        \item After all slips are collected, review several challenging questions as a class.
        \item Discuss common reasoning errors and strategies used by top-scoring pairs.
        \item Reinforce conceptual links between stations (e.g., gas behavior and intermolecular forces).
    \end{itemize}
\end{enumerate}
\\
\hline
\end{tabular}
\newpage
\noindent{\large\textbf{Content (fully worked out problems \& answers)}}\\[0.2em]
\vspace{0.2cm}
\noindent\begin{tabular}{|p{0.9\linewidth}|}
\hline
\textbf{(1)} Which solution has the highest [OH$^-$]? \\[0.2cm]
\begin{minipage}{0.85\linewidth}
\begin{enumerate}[label=\alph*.]
\item pH = 6.10 \item pH = 8.23 \item pH = 9.50 \item pH = 3.33
\end{enumerate}
\end{minipage}\\[0.2cm]
\hline
\end{tabular}

\noindent\begin{tabular}{|p{0.9\linewidth}|}
\hline
\textbf{(2)} Which substance is a strong base? \\[0.2cm]
\begin{minipage}[t]{0.85\linewidth}
\begin{enumerate}[label=\alph*.]
\item NH$_3$ \item Ca(OH)$_2$ \item HF \item HNO$_2$
\end{enumerate}
\end{minipage}\\[0.2cm]
\hline
\end{tabular}

\noindent\begin{tabular}{|p{0.9\linewidth}|}
\hline
\textbf{(3)} Which of the following statements about salt hydrolysis is correct? \\[0.2cm]
\begin{minipage}{0.85\linewidth}
\begin{enumerate}[label=\alph*.]
\item Only salts of strong acids hydrolyze
\item Only salts of strong bases hydrolyze
\item Salts of weak acids or weak bases hydrolyze
\item No salts hydrolyze in water
\end{enumerate}
\end{minipage}\\[0.2cm]
\hline
\end{tabular}

\noindent\begin{tabular}{|p{0.9\linewidth}|}
\hline
\textbf{(4)} Which salt forms a basic solution upon hydrolysis? \\[0.2cm]
\begin{minipage}{0.85\linewidth}
\begin{enumerate}[label=\alph*.]
\item KClO \item NH$_4$Br \item NaNO$_3$ \item LiCl
\end{enumerate}
\end{minipage}\\[0.2cm]
\hline
\end{tabular}
\newpage

\noindent\begin{tabular}{|p{0.9\linewidth}|}
\hline
\textbf{(5)} Which describes a neutral solution? \\[0.2cm]
\begin{minipage}{0.85\linewidth}
\begin{enumerate}[label=\alph*.]
\item $[H_3O^+] > [OH^-]$
\item $[H_3O^+] < [OH^-]$
\item $[H_3O^+] = [OH^-]$
\item $[H_3O^+] = 0$
\end{enumerate}
\end{minipage}\\[0.2cm]
\hline
\end{tabular}

\noindent\begin{tabular}{|p{0.9\linewidth}|}
\hline
\textbf{(6)} A solution has [$H_3O^+$] = $3.20\times10^{-6}$ M. Its pH is closest to: \\[0.2cm]
\begin{minipage}{0.85\linewidth}
\begin{enumerate}[label=\alph*.]
\item 4.49 \item 5.49 \item 6.49 \item 7.49
\end{enumerate}
\end{minipage}\\[0.2cm]
\hline
\end{tabular}

\noindent\begin{tabular}{|p{0.9\linewidth}|}
\hline
\textbf{(7)} Which salt will produce a basic solution upon dissolution due to its parent species? \\[0.2cm]
\begin{minipage}{0.85\linewidth}
\begin{enumerate}[label=\alph*.]
\item NaF \item NH$_4$Br \item KCl \item Ca(NO$_3$)$_2$
\end{enumerate}
\end{minipage}\\[0.2cm]
\hline
\end{tabular}

\noindent\begin{tabular}{|p{0.9\linewidth}|}
\hline
\textbf{(8)} In the reaction HCOOH + H$_2$O $\rightleftharpoons$ HCOO$^-$ + H$_3$O$^+$, HCOOH is acting as a(n): \\[0.2cm]
\begin{minipage}{0.85\linewidth}
\begin{enumerate}[label=\alph*.]
\item Base \item Acid \item Salt \item Neutral compound
\end{enumerate}
\end{minipage}\\[0.2cm]
\hline
\end{tabular}
\newpage

\noindent\begin{tabular}{|p{0.9\linewidth}|}
\hline
\textbf{(9)} Which of the following solutions is most basic? \\[0.2cm]
\begin{minipage}{0.85\linewidth}
\begin{enumerate}[label=\alph*.]
\item pH = 9.00 \item pH = 11.00 \item pH = 7.00 \item pH = 5.00
\end{enumerate}
\end{minipage}\\[0.2cm]
\hline
\end{tabular}

\noindent\begin{tabular}{|p{0.9\linewidth}|}
\hline
\textbf{(10)} Which of the following best explains why pKa values are useful? \\[0.2cm]
\begin{minipage}{0.85\linewidth}
\begin{enumerate}[label=\alph*.]
\item They show how fast an acid reacts
\item They show the strength of an acid using a log scale
\item They show solubility of an acid
\item They show the molecular mass of an acid
\end{enumerate}
\end{minipage}\\[0.2cm]
\hline
\end{tabular}

\noindent\begin{tabular}{|p{0.9\linewidth}|}
\hline
\textbf{(11)} Which term describes an acid that ionizes partially in water? \\[0.2cm]
\begin{minipage}[t]{0.85\linewidth}
\begin{enumerate}[label=\alph*.]
\item Strong acid \item Weak acid \item Strong base \item Neutral compound
\end{enumerate}
\end{minipage}\\[0.2cm]
\hline
\end{tabular}

\noindent\begin{tabular}{|p{0.9\linewidth}|}
\hline
\textbf{(12)} Which description correctly identifies the parent acid of the salt sodium acetate (NaOAc)? \\[0.2cm]
\begin{minipage}{0.85\linewidth}
\begin{enumerate}[label=\alph*.]
\item HCl \item CH$_3$COOH \item HNO$_3$ \item HF
\end{enumerate}
\end{minipage}\\[0.2cm]
\hline
\end{tabular}
\newpage

\noindent\begin{tabular}{|p{0.9\linewidth}|}
\hline
\textbf{(13)} Which species is the conjugate acid of NH$_3$? \\[0.2cm]
\begin{minipage}{0.85\linewidth}
\begin{enumerate}[label=\alph*.]
\item NH$_2^-$ \item NH$_4^+$ \item N$^{3-}$ \item NH$_2$OH
\end{enumerate}
\end{minipage}\\[0.2cm]
\hline
\end{tabular}

\noindent\begin{tabular}{|p{0.9\linewidth}|}
\hline
\textbf{(14)} In the stepwise dissociation of H$_2$SO$_4$, which statement is correct? \\[0.2cm]
\begin{minipage}{0.85\linewidth}
\begin{enumerate}[label=\alph*.]
\item Both protons dissociate completely \item The first proton dissociates strongly, the second weakly \item Neither proton dissociates \item The second proton dissociates more strongly than the first
\end{enumerate}
\end{minipage}\\[0.2cm]
\hline
\end{tabular}

\noindent\begin{tabular}{|p{0.9\linewidth}|}
\hline
\textbf{(15)} For the equilibrium $\text{HA} + \text{H}_2\text{O} \rightleftharpoons \text{H}_3\text{O}^+ + \text{A}^-$, which expression represents $K_a$? \\[0.2cm]
\begin{minipage}[t]{0.85\linewidth}
\begin{enumerate}[label=\alph*.]
\item $\dfrac{[\text{HA}]}{[\text{H}_3\text{O}^+][\text{A}^-]}$ \item $\dfrac{[\text{H}_3\text{O}^+][\text{A}^-]}{[\text{HA}]}$ \item $[\text{H}_3\text{O}^+][\text{A}^-][\text{HA}]$ \item $\dfrac{1}{[\text{HA}]}$
\end{enumerate}
\end{minipage}\\[0.2cm]
\hline
\end{tabular}

\noindent\begin{tabular}{|p{0.9\linewidth}|}
\hline
\textbf{(16)} In the reaction $\text{NH}_3 + \text{H}_2\text{O} \rightarrow \text{NH}_4^+ + \text{OH}^-$, which species acts as the Brönsted-Lowry base?\\[0.2cm]
\begin{minipage}[t]{0.85\linewidth}
\begin{enumerate}[label=\alph*.]
\item NH$_3$
\item H$_2$O
\item NH$_4^+$
\item OH$^-$
\end{enumerate}
\end{minipage}\\[0.2cm]
\hline
\end{tabular}
\newpage

\noindent\begin{tabular}{|p{0.9\linewidth}|}
\hline
\textbf{(17)} Which species is amphiprotic? \\[0.2cm]
\begin{minipage}{0.85\linewidth}
\begin{enumerate}[label=\alph*.]
\item OH$^-$ \item HSO$_4^-$ \item Cl$^-$ \item NH$_4^+$
\end{enumerate}
\end{minipage}\\[0.2cm]
\hline
\end{tabular}

\noindent\begin{tabular}{|p{0.9\linewidth}|}
\hline
\textbf{(18)} A solution has pH = 5.00. What is the hydronium concentration? \\[0.2cm]
\begin{minipage}{0.85\linewidth}
\begin{enumerate}[label=\alph*.]
\item $1.00\times10^{-3}$ M \item $1.00\times10^{-5}$ M \item $1.00\times10^{-7}$ M \item $1.00\times10^{-9}$ M
\end{enumerate}
\end{minipage}\\[0.2cm]
\hline
\end{tabular}

\noindent\begin{tabular}{|p{0.9\linewidth}|}
\hline
\textbf{(19)} If a salt's anion is the conjugate base of a weak acid, what will be the pH of its aqueous solution? \\[0.2cm]
\begin{minipage}{0.85\linewidth}
\begin{enumerate}[label=\alph*.]
\item Acidic \item Basic \item Neutral \item Impossible to determine
\end{enumerate}
\end{minipage}\\[0.2cm]
\hline
\end{tabular}

\noindent\begin{tabular}{|p{0.9\linewidth}|}
\hline
\textbf{(20)} A weak acid has Ka = $3.00\times10^{-7}$. Which pKa is closest? \\[0.2cm]
\begin{minipage}{0.85\linewidth}
\begin{enumerate}[label=\alph*.]
\item 3.96 \item 4.22 \item 5.44 \item 6.80
\end{enumerate}
\end{minipage}\\[0.2cm]
\hline
\end{tabular}
\newpage

\noindent\begin{tabular}{|p{0.9\linewidth}|}
\hline
\textbf{(21)} A solution has a hydronium concentration of $1.0\times10^{-4}\ \text{M}$. What is its pH? \\[0.2cm]
\begin{minipage}[t]{0.85\linewidth}
\begin{enumerate}[label=\alph*.]
\item 2 \item 3 \item 4 \item 5
\end{enumerate}
\end{minipage}\\[0.2cm]
\hline
\end{tabular}

\noindent\begin{tabular}{|p{0.9\linewidth}|}
\hline
\textbf{(22)} Which statement correctly distinguishes dissociation from ionization? \\[0.2cm]
\begin{minipage}{0.85\linewidth}
\begin{enumerate}[label=\alph*.]
\item Dissociation involves breaking covalent bonds to form new ions, while ionization separates pre-existing ions in an ionic lattice.
\item Dissociation separates pre-existing ions from an ionic compound without altering their identity or overall charge, while ionization forms ions from a molecular substance through a chemical change.
\item Dissociation always produces H$_3$O$^+$, while ionization never does.
\item Dissociation only occurs in strong acids, while ionization only occurs in weak acids.
\end{enumerate}
\end{minipage}\\[0.2cm]
\hline
\end{tabular}

\noindent\begin{tabular}{|p{0.9\linewidth}|}
\hline
\textbf{(23)} A weak base has $[\text{OH}^-] = 1.00\times10^{-5}\ \text{M}$. What is the pOH? \\[0.2cm]
\begin{minipage}{0.85\linewidth}
\begin{enumerate}[label=\alph*.]
\item 3.00 \item 5.00 \item 7.00 \item 9.00
\end{enumerate}
\end{minipage}\\[0.2cm]
\hline
\end{tabular}

\noindent\begin{tabular}{|p{0.9\linewidth}|}
\hline
\textbf{(24)} For the equilibrium HF + H$_2$O $\rightleftharpoons$ H$_3$O$^+$ + F$^-$, which species is the conjugate base? \\[0.2cm]
\begin{minipage}{0.85\linewidth}
\begin{enumerate}[label=\alph*.]
\item HF \item H$_3$O$^+$ \item F$^-$ \item H$_2$O
\end{enumerate}
\end{minipage}\\[0.2cm]
\hline
\end{tabular}
\newpage

\noindent\begin{tabular}{|p{0.9\linewidth}|}
\hline
\textbf{(25)} A base that ionizes only slightly in water is classified as a: \\[0.2cm]
\begin{minipage}{0.85\linewidth}
\begin{enumerate}[label=\alph*.]
\item Strong base \item Weak base \item Neutral molecule \item Strong acid
\end{enumerate}
\end{minipage}\\[0.2cm]
\hline
\end{tabular}

\noindent\begin{tabular}{|p{0.9\linewidth}|}
\hline
\textbf{(26)} The parent acid of NO$_2^-$ is: \\[0.2cm]
\begin{minipage}{0.85\linewidth}
\begin{enumerate}[label=\alph*.]
\item HNO$_3$ \item HNO$_2$ \item HNO \item HNO$_4$
\end{enumerate}
\end{minipage}\\[0.2cm]
\hline
\end{tabular}

\noindent\begin{tabular}{|p{0.9\linewidth}|}
\hline
\textbf{(27)} A solution has pOH = 3.60. Determine the pH, [H$_3$O$^+$], and [OH$^-$] of the solution.  Ensure concentrations are in molarity and pH/pOH to two decimals. \\[4cm]
\hline
\end{tabular}

\noindent\begin{tabular}{|p{0.9\linewidth}|}
\hline
\textbf{(28)} Calcium hydroxide reacts with hydrochloric acid to form calcium chloride and water. Which statement best describes the acid-base behavior occurring in this reaction?\\[0.2cm]
Ca(OH)$_2$ + 2HCl $\rightarrow$ CaCl$_2$ + 2H$_2$O\\[0.2cm]
\begin{minipage}[t]{0.85\linewidth}
\begin{enumerate}[label=\alph*.]
\item HCl donates protons to OH$^-$ ions, forming water as Ca$^{2+}$ pairs with Cl$^-$
\item Ca(OH)$_2$ donates protons to HCl, forming hydrogen gas
\item Both reactants accept protons, preventing water formation
\item Ca$^{2+}$ acts as the proton donor while Cl$^-$ acts as the proton acceptor
\end{enumerate}
\end{minipage}\\[0.2cm]
\hline
\end{tabular}
\newpage

\noindent\begin{tabular}{|p{0.9\linewidth}|}
\hline
\textbf{(29)} Which statement correctly describes a strong base? \\[0.2cm]
\begin{minipage}{0.85\linewidth}
\begin{enumerate}[label=\alph*.]
\item Completely dissociates in water \item Partially dissociates in water \item Does not react with water \item Cannot accept protons
\end{enumerate}
\end{minipage}\\[0.2cm]
\hline
\end{tabular}

\noindent\begin{tabular}{|p{0.9\linewidth}|}
\hline
\textbf{(30)} A solution with pOH = 3.00 has what [OH$^-$]? \\[0.2cm]
\begin{minipage}{0.85\linewidth}
\begin{enumerate}[label=\alph*.]
\item $1.00\times10^{-3}$ M \item $1.00\times10^{-7}$ M \item $1.00\times10^{-11}$ M \item $1.00\times10^{-14}$ M
\end{enumerate}
\end{minipage}\\[0.2cm]
\hline
\end{tabular}

\noindent\begin{tabular}{|p{0.9\linewidth}|}
\hline
\textbf{(31)} Which statement best describes the logarithmic nature of pKa? \\[0.2cm]
\begin{minipage}{0.85\linewidth}
\begin{enumerate}[label=\alph*.]
\item A decrease of 1 pKa means the acid is ten times stronger
\item A decrease of 1 pKa means the acid is half as strong
\item A decrease of 1 pKa means no measurable change
\item A decrease of 1 pKa means the acid becomes more basic
\end{enumerate}
\end{minipage}\\[0.2cm]
\hline
\end{tabular}

\noindent\begin{tabular}{|p{0.9\linewidth}|}
\hline
\textbf{(32)} Which of the following is a triprotic acid? \\[0.2cm]
\begin{minipage}{0.85\linewidth}
\begin{enumerate}[label=\alph*.]
\item HCl \item H$_2$CO$_3$ \item H$_3$PO$_4$ \item HF
\end{enumerate}
\end{minipage}\\[0.2cm]
\hline
\end{tabular}
\newpage

\noindent\begin{tabular}{|p{0.9\linewidth}|}
\hline
\textbf{(33)} Which is the conjugate acid of CO$_3^{2-}$? \\[0.2cm]
\begin{minipage}{0.85\linewidth}
\begin{enumerate}[label=\alph*.]
\item CO$_2$ \item HCO$_3^-$ \item H$_2$CO$_3$ \item CO$_3$H$_2$
\end{enumerate}
\end{minipage}\\[0.2cm]
\hline
\end{tabular}

\noindent\begin{tabular}{|p{0.9\linewidth}|}
\hline
\textbf{(34)} Which salt will undergo acidic hydrolysis when dissolved in water? \\[0.2cm]
\begin{minipage}{0.85\linewidth}
\begin{enumerate}[label=\alph*.]
\item KNO$_3$ \item NH$_4$Cl \item NaClO$_3$ \item CaCl$_2$
\end{enumerate}
\end{minipage}\\[0.2cm]
\hline
\end{tabular}

\noindent\begin{tabular}{|p{0.9\linewidth}|}
\hline
\textbf{(35)} A solution has pH = 11.00 What is the hydroxide concentration? \\[0.2cm]
\begin{minipage}[t]{0.85\linewidth}
\begin{enumerate}[label=\alph*.]
\item $1.00\times10^{-3}$ M \item $1.00\times10^{-11}$ M \item $1.00\times10^{-7}$ M \item $1.00\times10^{-14}$ M
\end{enumerate}
\end{minipage}\\[0.2cm]
\hline
\end{tabular}

\noindent\begin{tabular}{|p{0.9\linewidth}|}
\hline
\textbf{(36)} A Brönsted-Lowry acid is defined as a substance that: \\[0.2cm]
\begin{minipage}[t]{0.85\linewidth}
\begin{enumerate}[label=\alph*.]
\item Donates a proton
\item Accepts a proton
\item Releases electrons
\item Accepts electrons
\end{enumerate}
\end{minipage}\\[0.2cm]
\hline
\end{tabular}
\newpage

\noindent\begin{tabular}{|p{0.9\linewidth}|}
\hline
\textbf{(37)} A weak acid has pKa = 4.25 and [HA] = 0.10 M. If [H$_3$O$^+$] = $5.6\times10^{-5}$ M, calculate pH, pOH, and [OH$^-$].  Ensure concentrations are in molarity and pH/pOH to two decimals. \\[4cm]
\hline
\end{tabular}

\noindent\begin{tabular}{|p{0.9\linewidth}|}
\hline
\textbf{(38)} In the reaction $\text{H}_2\text{PO}_4^-$ + H$_2$O \rightleftharpoons \text{HPO}_4^{2-} + \text{H}_3O^+$, the conjugate base is: \\[0.2cm]
\begin{minipage}[t]{0.85\linewidth}
\begin{enumerate}[label=\alph*.]
\item H$_2$PO$_4^-$
\item HPO$_4^{2-}$
\item H$_3$O$^+$
\item None of the above
\end{enumerate}
\end{minipage}\\[0.2cm]
\hline
\end{tabular}

\noindent\begin{tabular}{|p{0.9\linewidth}|}
\hline
\textbf{(39)} Which of the following best defines ionization for a weak acid in water? \\[0.2cm]
\begin{minipage}[t]{0.85\linewidth}
\begin{enumerate}[label=\alph*.]
\item Complete separation into ions \item Partial separation into ions \item Reaction with a base to form a salt \item Combination with water molecules without forming ions
\end{enumerate}
\end{minipage}\\[0.2cm]
\hline
\end{tabular}

\noindent\begin{tabular}{|p{0.9\linewidth}|}
\hline
\textbf{(40)} A weak acid has pKa = 3.97. Its solution has [H$_3$O$^+$] = $8.00\times10^{-5}$ M. Determine pH, pOH, and [OH$^-$].  Ensure concentrations are in molarity and pH/pOH to two decimals.\\[4cm]
\hline
\end{tabular}
\newpage

\noindent\begin{tabular}{|p{0.9\linewidth}|}
\hline
\textbf{(41)} Which pair represents a conjugate acid-base pair?\\[0.2cm]
\begin{minipage}[t]{0.85\linewidth}
\begin{enumerate}[label=\alph*.]
\item NH$_3$/NH$_4^+$
\item OH$^-$/O$^{2-}$
\item Na$^+$/Na
\item Cl$^-$/Cl$_2$
\end{enumerate}
\end{minipage}\\[0.2cm]
\hline
\end{tabular}

\noindent\begin{tabular}{|p{0.9\linewidth}|}
\hline
\textbf{(42)} Which describes the relationship between Ka and acid strength? \\[0.2cm]
\begin{minipage}{0.85\linewidth}
\begin{enumerate}[label=\alph*.]
\item Larger Ka means stronger acid \item Smaller Ka means stronger acid \item Ka does not relate to strength \item Ka only applies to bases
\end{enumerate}
\end{minipage}\\[0.2cm]
\hline
\end{tabular}

\noindent\begin{tabular}{|p{0.9\linewidth}|}
\hline
\textbf{(43)} If the pH of a solution increases from 8.00 to 10.00, the [H$^+$] changes by a factor of: \\[0.2cm]
\begin{minipage}{0.85\linewidth}
\begin{enumerate}[label=\alph*.]
\item 2x \item 10x \item 50x \item 100x
\end{enumerate}
\end{minipage}\\[0.2cm]
\hline
\end{tabular}

\noindent\begin{tabular}{|p{0.9\linewidth}|}
\hline
\textbf{(44)} A strong base solution has [OH$^-$] = $6.08\times10^{-4}$ M. Calculate pOH, pH, and [H$_3$O$^+$] to two decimal places, with concentrations in molarity. \\[4cm]
\hline
\end{tabular}
\newpage

\noindent\begin{tabular}{|p{0.9\linewidth}|}
\hline
\textbf{(45)} A solution with pH = 2.00 has a hydronium concentration that is how many times larger than a solution with pH = 4.00? \\[0.2cm]
\begin{minipage}{0.85\linewidth}
\begin{enumerate}[label=\alph*.]
\item 2x \item 10x \item 100x \item 1000x
\end{enumerate}
\end{minipage}\\[0.2cm]
\hline
\end{tabular}

\noindent\begin{tabular}{|p{0.9\linewidth}|}
\hline
\textbf{(46)} If an acid has $K_a = 1.00\times10^{-5}$, which of the following is closest to its pKa? \\[0.2cm]
\begin{minipage}{0.85\linewidth}
\begin{enumerate}[label=\alph*.]
\item 2.20 \item 3.50 \item 5.33 \item 7.25
\end{enumerate}
\end{minipage}\\[0.2cm]
\hline
\end{tabular}

\noindent\begin{tabular}{|p{0.9\linewidth}|}
\hline
\textbf{(47)} A solution has pH = 5.83. Find [H$_3$O$^+$], pOH, and [OH$^-$]. Concentrations in M, pH/pOH to two decimal places. \\[4cm]
\hline
\end{tabular}

\noindent\begin{tabular}{|p{0.9\linewidth}|}
\hline
\textbf{(48)} Which statement correctly defines ionization for a base in water? \\[0.2cm]
\begin{minipage}[t]{0.85\linewidth}
\begin{enumerate}[label=\alph*.]
\item Forming OH$^-$ by reacting chemically with water rather than releasing pre-existing ions
\item Separating into pre-existing ions without chemical change
\item Producing H$_3$O$^+$ by transferring a proton to water
\item Splitting into neutral molecules without forming ions
\end{enumerate}
\end{minipage}\\[0.2cm]
\hline
\end{tabular}
\newpage

\noindent\begin{tabular}{|p{0.9\linewidth}|}
\hline
\textbf{(49)} Which species will act as a base according to Brönsted-Lowry? \\[0.2cm]
\begin{minipage}{0.85\linewidth}
\begin{enumerate}[label=\alph*.]
\item H$_3$O$^+$ \item NH$_3$ \item HCl \item HNO$_3$
\end{enumerate}
\end{minipage}\\[0.2cm]
\hline
\end{tabular}

\noindent\begin{tabular}{|p{0.9\linewidth}|}
\hline
\textbf{(50)} A solution has [H$_3$O$^+$] = $2.22\times10^{-6}$ M. Calculate the pH, pOH, and [OH$^-$]. \\[4cm]
\hline
\end{tabular}

\noindent\begin{tabular}{|p{0.9\linewidth}|}
\hline
\textbf{(51)} Which of the following is a strong acid? \\[0.2cm]
\begin{minipage}[t]{0.85\linewidth}
\begin{enumerate}[label=\alph*.]
\item HF \item HCl \item HCN \item CH$_3$COOH
\end{enumerate}
\end{minipage}\\[0.2cm]
\hline
\end{tabular}

\noindent\begin{tabular}{|p{0.9\linewidth}|}
\hline
\textbf{(52)} A solution has [H$_3$O$^+$] = $4.58\times10^{-5}$ M. Calculate the pH, pOH, and [OH$^-$]. Express concentrations in molarity and pH/pOH to two decimal places. \\[4cm]
\hline
\end{tabular}
\newpage

\noindent\begin{tabular}{|p{0.9\linewidth}|}
\hline
\textbf{(53)} Which salt produces an acidic solution? \\[0.2cm]
\begin{minipage}{0.85\linewidth}
\begin{enumerate}[label=\alph*.]
\item NaCl \item KCN \item NH$_4$NO$_3$ \item NaOAc
\end{enumerate}
\end{minipage}\\[0.2cm]
\hline
\end{tabular}

\noindent\begin{tabular}{|p{0.9\linewidth}|}
\hline
\textbf{(54)} Which salt is formed from a strong acid and strong base and therefore produces a neutral solution? \\[0.2cm]
\begin{minipage}{0.85\linewidth}
\begin{enumerate}[label=\alph*.]
\item NaCl \item NH$_4$Cl \item NaOAc \item KCN
\end{enumerate}
\end{minipage}\\[0.2cm]
\hline
\end{tabular}

\noindent\begin{tabular}{|p{0.9\linewidth}|}
\hline
\textbf{(55)} When hydrochloric acid reacts with sodium hydroxide to form sodium chloride and water, which statement best describes the acid-base process taking place?\\[0.2cm]
HCl + NaOH $\rightarrow$ NaCl + H$_2$O\\[0.2cm]
\begin{minipage}[t]{0.85\linewidth}
\begin{enumerate}[label=\alph*.]
\item HCl donates a proton to OH$^-$, forming water as Na$^+$ pairs with Cl$^-$
\item NaOH donates protons to HCl, forming hydrogen gas
\item Both reactants act as proton donors, preventing neutralization
\item Cl$^-$ acts as the proton donor while Na$^+$ accepts the proton
\end{enumerate}
\end{minipage}\\[0.2cm]
\hline
\end{tabular}

\noindent\begin{tabular}{|p{0.9\linewidth}|}
\hline
\textbf{(56)} A weak acid HA has a pKa of 4.70. Which of the following pH values will favor the deprotonated form, A$^-$? \\[0.2cm]
\begin{minipage}{0.85\linewidth}
\begin{enumerate}[label=\alph*.]
\item 2.00 \item 3.00 \item 4.70 \item 7.00
\end{enumerate}
\end{minipage}\\[0.2cm]
\hline
\end{tabular}
\newpage

\noindent\begin{tabular}{|p{0.9\linewidth}|}
\hline
\textbf{(57)} According to the Arrhenius definition, which statement correctly describes an acid and a base?\\[0.2cm]
\begin{minipage}[t]{0.85\linewidth}
\begin{enumerate}[label=\alph*.]
\item An acid produces H$^+$ in water, and a base produces OH$^-$ in water
\item An acid donates a proton, and a base accepts a proton
\item An acid accepts an electron pair, and a base donates an electron pair
\item An acid and base both increase the concentration of ions equally
\end{enumerate}
\end{minipage}\\[0.2cm]
\hline
\end{tabular}

\noindent\begin{tabular}{|p{0.9\linewidth}|}
\hline
\textbf{(58)} In the equilibrium $\text{HF} + \text{H}_2\text{O} \rightleftharpoons \text{H}_3\text{O}^+ + \text{F}^-$, which is the conjugate acid of F$^-$? \\[0.2cm]
\begin{minipage}[t]{0.85\linewidth}
\begin{enumerate}[label=\alph*.]
\item HF
\item H$_3$O$^+$
\item H$_2$O
\item F$^-$ has no conjugate acid
\end{enumerate}
\end{minipage}\\[0.2cm]
\hline
\end{tabular}

\noindent\begin{tabular}{|p{0.9\linewidth}|}
\hline
\textbf{(59)} Which acid has the strongest conjugate base? \\[0.2cm]
\begin{minipage}{0.85\linewidth}
\begin{enumerate}[label=\alph*.]
\item HCl \item HNO$_3$ \item HF \item CH$_3$COOH
\end{enumerate}
\end{minipage}\\[0.2cm]
\hline
\end{tabular}

\noindent\begin{tabular}{|p{0.9\linewidth}|}
\hline
\textbf{(60)} Which statement is true about conjugate acid-base pairs? \\[0.2cm]
\begin{minipage}{0.85\linewidth}
\begin{enumerate}[label=\alph*.]
\item They differ by two protons \item They differ by one proton \item They differ by one electron \item They differ by two electrons
\end{enumerate}
\end{minipage}\\[0.2cm]
\hline
\end{tabular}


\newpage
\noindent
{\large\textbf{Checking for Understanding (questions)}}\\[0.2em]
\noindent
\begin{tabular}{|p{0.9\linewidth}|}
\hline
\begin{enumerate}
    \item How do you distinguish between dissociation and ionization in aqueous solutions? Give examples of each.  
    \item How is the hydronium concentration \([H_3O^+]\) related to pH? How do you calculate one from the other?  
    \item How do you determine pOH and \([OH^-]\) from a known pH?  
    \item What is the relationship between Ka and acid strength? How can you estimate which species is more acidic?  
    \item How do you calculate percent ionization for a weak acid given \([H_3O^+]\) and initial concentration?  
    \item How do pKa values inform which form (protonated or deprotonated) predominates at a given pH?  
    \item How do you identify the conjugate acid or base in a given reaction?  
    \item How do you predict the effect of a strong acid vs. weak acid on pH?  
    \item How does a diprotic or triprotic acid behave differently in terms of sequential ionization?  
    \item How do you determine whether a salt solution will be acidic, basic, or neutral based on the strengths of its parent acid and base?  
    \item What is salt hydrolysis, and how does it affect the pH of the solution?  
    \item How do you check that your pH, pOH, \([H_3O^+]\), and \([OH^-]\) calculations are consistent?  
    \item Why is it important to include units for concentration (M) and round pH/pOH to two decimal places?  
    \item How do strong acids and strong bases differ from weak acids and weak bases in terms of dissociation/ionization?  
    \item How can you predict the direction of proton transfer in a Brønsted-Lowry acid-base reaction?  
\end{enumerate} \\
\hline  
\end{tabular}

\newpage
\noindent
\section*{\underline{Closer}}
\begin{tabular}{|p{4.2cm}|p{9.8cm}|}
\hline
\textbf{Collaborative Learning Techniques} & Group Survey\\
\hline
\textbf{Learning Strategy} & Brain dump\\
\hline
\textbf{Time} & 12:30---12:32\\
\hline
\textbf{Materials} & Index cards, writing utensils\\
\hline
\textbf{Lecture/Textbook/Old Plans/Slide References} & Chapters 1-10 \\
\hline
\end{tabular}

\vspace{0.5cm}
\noindent
\begin{tabular}{|p{0.9\linewidth}|}
\hline
{\large\textbf{Definitions}}
\begin{enumerate}
   \item\textbf{Group survey:} Each group member is surveyed to discover their position on an issue, problem or topic. This process ensures that each member of the group is allowed to offer or state their point of view.
   \item\textbf{Brain dump:} This is a great closer for students to sum up their knowledge on the topic(s) from the session. Have the students take out a blank piece of paper /Google Doc and write/type("dump") everything they know relevant from the discussion that day. The SI Leader can look at their work to assess the students' understanding of the course material and have them compare it to each other. It is important to at least discuss with the students so the SI Leader knows the students' comfort levels with the material. 
\end{enumerate}
\\ \hline
\end{tabular}

\newpage
\noindent
{\large\textbf{Instructions (please provide a \hl{DETAILED} activity breakdown below (and/or attached}}

\begin{tabular}{|p{0.9\linewidth}|}
\hline
\begin{enumerate}[itemsep=0.3em]
   \item Students are put into small groups and are asked to write down their confidence on a single notecard without names.
   \item The cards are evidence of the group's confidence in the content.
\end{enumerate}
\\ \hline
\end{tabular}

\vspace{0.5cm}
\noindent
{\large\textbf{Content (fully worked out problems \& answers)}

\begin{tabular}{|p{0.9\linewidth}|}
\hline
\begin{itemize}
   \item No new worked problems; brain dump.
\end{itemize}
\\ \hline
\end{tabular}

\vspace{0.5cm}
\noindent
{\large\textbf{Checking for Understanding (questions \& answers) }

\begin{tabular}{|p{0.9\linewidth}|}
\hline
\begin{itemize}
   \item Review confidence levels
   \item Interpret the variety of confidence and their distribution in groups.
   \item Get feedback for what they like and don't like
   \item Check in if the time is suitable for the students
\end{itemize}
\\ \hline
\end{tabular}
\newpage
\section*{Plan Checklist}

\begin{tabular}{|p{0.9\linewidth}|}
\hline
\begin{itemize}
\item[$\Box$] Variety of Collaborative Learning Techniques and Learning Strategies
\item[$\Box$] Chapter/slide\# where information is coming from
\item[$\Box$] Included timestamps (and buffer time as needed)
\item[$\Box$] Details on how leader will break up students
\item[$\Box$] Checking for Understanding questions
\item[$\Box$] Fully worked out problems and examples for each activity
\item[$\Box$] Necessary links required for online implementation (if applicable)
\item[$\Box$] Source material correctly cited (if applicable)
\item[$\Box$] Plan is uploaded to Teams
\end{itemize}
\\ \hline
\end{tabular}
\newpage

\end{document}